\section{Conclusion and Future Work}
\label{sec:conclusion}

In this paper, we presented a rigorous investigation into the post-training quantization (PTQ) robustness of test-time adaptive model merging. We identified and characterized a fundamental vulnerability in existing unregularized adaptive merging frameworks: \textbf{Quantization-Operator Overfitting}. While learning layer-wise blending coefficients via test-time adaptation yields high multi-task accuracy in floating-point (FP32) space, it forces the network parameters into sharp local minima that are highly fragile under post-training low-precision quantization.

To resolve this limitation, we introduced \textbf{Hessian-Regularized Coefficient Optimization (HessMerge)}, a mathematically grounded framework that guarantees PTQ robustness. By modeling post-training quantization as a parameter-space perturbation, we showed via a second-order Taylor expansion that the resulting performance drop is directly governed by the curvature of the loss landscape with respect to the merging coefficients. HessMerge explicitly incorporates a curvature penalty based on the trace of the coefficient-space Hessian ($\text{Tr}(\mathcal{H}_{\Theta})$) during test-time adaptation. This second-order penalty flattens the local loss minimum, ensuring that weight-space rounding noise from any downstream quantization operator leads to negligible performance degradation.

Our comprehensive multi-schema evaluation using a Vision Transformer backbone on four distinct visual domain datasets (MNIST, FashionMNIST, CIFAR-10, SVHN) demonstrated that HessMerge consistently outperforms established model merging baselines (such as Task Arithmetic, AdaMerging, RegCalMerge, Q-Merge, and PolyMerge) across multiple deployment quantization formats. Most notably, HessMerge retains high joint multi-task performance even under aggressive INT4 symmetric per-channel quantization, where traditional adaptive schemes suffer from catastrophic representation collapse.

There are several exciting directions for future research. First, we plan to scale HessMerge to large language models (LLMs) and large-scale multimodal vision-language models (e.g., CLIP, LLaVA), where edge-deployment compression is even more critical. Second, we aim to extend our Hessian trace regularization to online, dynamic quantization settings, where hardware precision shifts adaptively during runtime. Finally, we hope to explore theoretical connections between our coefficient-space flat-minimum regularization and general out-of-distribution generalization in weight-space model merging.
